%==================================================
%----- 2.8 Interfaces/portas de acoplamento
%==================================================
\section{Interfaces/portas de acoplamento}

\subsection{Referenciais e cinemática relativa}
São empregados três referenciais:
a base $\{B\}$, o efetuador $\{E\}$ e a porta de acoplamento (\emph{docking}) $\{D\}$. 
A posição do ponto de acoplamento do efetuador expressa em
$\{D\}$ é representada por
${}^{D}\!\vec p_{E}(\vec q,t)\in\mathbb R^{3}$,
e a orientação relativa pela rotação
${}^{D}\!\mathbf R_{E}(\vec q,t)\in\mathrm{SO}(3)$.

\subsection{Alinhamento}
O alinhamento geométrico é modelado como restrição holonômica
$\Phi(\vec q,t)=\vec 0$,
em conformidade com \eqref{eq:PhiJ1} e \eqref{eq:PhiJ2}. 
O erro de posição em $\{D\}$ é dado por:
\begin{equation}
\Phi_{p}(\vec q,t)\;
\triangleq\;
{}^{D}\!\vec p_{E}(\vec q,t)-\hat{\vec p}_{D}(t)\in\mathbb R^{3},
\label{eq:phi_pos}
\end{equation}
com $\hat{\vec p}_{D}(t)$ no centro da porta. 
Para a orientação, utiliza-se o quaternion relativo
${}^{D}\!q_{E}=[\,\eta\;\;\vec\epsilon^{\top}\,]^{\top}$
associado a ${}^{D}\!\mathbf R_{E}$;
além disso, adotou-se como erro mínimo a sua parte vetorial:
\begin{equation}
\Phi_{o}(\vec q,t)\;\triangleq\;\vec\epsilon\in\mathbb R^{3}.
\label{eq:phi_ori}
\end{equation}
A restrições é definida por:
\begin{equation}
\Phi(\vec q,t)\;\triangleq\;\begin{bmatrix}\Phi_{p}\\ \Phi_{o}\end{bmatrix}=\vec 0,
\qquad 
\mathbf J(\vec q,t)=\frac{\partial \Phi}{\partial \vec q},
\label{eq:phi_stack_dock}
\end{equation}

cujos níveis de velocidade e aceleração seguem de
\eqref{eq:vellevel} e \eqref{eq:acclevel}:
\begin{equation}
    \label{eq:vel_acel_porta_docking}
\dot\Phi
=
\mathbf J\,\dot{\vec q}
+\frac{\partial \Phi}{\partial t},
\qquad
\ddot\Phi
=
\mathbf J\,\ddot{\vec q}
+\dot{\mathbf J}\,\dot{\vec q}
+\frac{\partial^{2}\Phi}{\partial t^{2}}
\end{equation}

\subsection{Condição de não-penetração}
A não penetração é descrita por uma desigualdade escalar.
Com $\vec n_{D}$ indicando a normal unitária do plano da $boca$ da porta e
${}^{D}\!\vec r_{E}(\vec q,t)$
o vetor do ponto do efetuador até esse plano, sendo expresso como:
\begin{equation}
g_{n}(\vec q,t)\;
\triangleq\;
\vec n_{D}^{\top}\,{}^{D}\!\vec r_{E}(\vec q,t)\;\ge 0.
\label{eq:gap}
\end{equation}
No contato perfeito, $g_{n}=0$;
após a captura, o acoplamento rígido é representado por
\eqref{eq:phi_stack_dock} como vínculo ativo.

\subsection{Jacobianas de acoplamento}
As linhas de $\mathbf J$ associadas a $\Phi_{p}$
coincidem com a Jacobiana de velocidade linear do efetuador expressa em $\{D\}$. 
As linhas relativas a $\Phi_{o}$ são proporcionais à Jacobiana de velocidade angular;
já a cinemática de quaternions fornece:
\begin{equation}
\dot{\vec\epsilon}
=
\tfrac12\big
(\eta\,\mathbf I+\lfloor\vec\epsilon\times\rfloor\big)\,
{}^{D}\!\vec\omega_{E/D},
\end{equation}
onde se obtém a relação com ${}^{D}\!\mathbf J_{\omega}$.

\subsection{Referências para o controle da base na aproximação}
Durante a aproximação, a atitude da base é estabilizada por um controlador PD que está sendo referenciado a $\{D\}$.
Com $\mathbf R_{b}$ e $\mathbf R_{b}^{\mathrm{ref}}(t)$,
sendo as orientações atual e de referência, os erros podem ser expressos por:
\begin{equation}
\vec e_{R}\;
\triangleq\;
\mathrm{vex}\!\Big
(\tfrac12\big(\mathbf R_{b}^{\top}\mathbf R_{b}^{\mathrm{ref}}-\mathbf R_{b}^{\mathrm{ref}\,\top}\mathbf R_{b}\big)\Big),
\qquad
\vec e_{\omega}\;
\triangleq\;
\vec\omega_{b}
-\vec\omega_{b}^{\mathrm{ref}},
\label{eq:err_base}
\end{equation}
e a lei:
\begin{equation}
\vec\tau_{b}
=
-\,\mathbf K_{R}\,\vec e_{R}
-\mathbf K_{\omega}\,\vec e_{\omega}.
\label{eq:pd_base}
\end{equation}
O manipulador robótico executa o alinhamento de $\{E\}$ com $\{D\}$ via \eqref{eq:phi_stack_dock};
caso se mantenham as restrições por estabilização, aplica-se \eqref{eq:acclevel_baumgarte}.

\subsection{Observações}
(i) A transição entre a aproximação e o contato, comuta o conjunto de vínculos ativos,
de \eqref{eq:gap} para \eqref{eq:phi_stack_dock}; 
(ii) No caso de perda de posto de $\mathbf J$, pode-se usar a pseudoinversa amortecida; 
(iii) As formulações de energia e dissipação previamente apresentadas permanecem válidas com as novas restrições.
